\documentclass[11pt,a4paper]{article}

\usepackage[T1]{fontenc}
\usepackage[utf8]{inputenc}
\usepackage[english]{babel}
\usepackage{lmodern}
\usepackage{microtype}
\usepackage{geometry}
\usepackage{graphicx}
\usepackage{booktabs}
\usepackage{amsmath}
\usepackage{amssymb}
\usepackage{hyperref}

\geometry{margin=2.5cm}

\title{Article title}
\author{First Last\\Institution\\\texttt{email@example.com}}
\date{\today}

\begin{document}

\maketitle

\begin{abstract}
This abstract summarizes the context, method, main results, and contribution of the article in a few sentences.
\end{abstract}

\section{Introduction}

Present the scientific problem, the literature context, and the main contribution of the article.

\section{Method}

Describe the data, assumptions, and method used.

\begin{equation}
  y = f(x) + \varepsilon
\end{equation}

\section{Results}

Present the results concisely.

\begin{table}[h]
  \centering
  \caption{Example results table.}
  \begin{tabular}{lcc}
    \toprule
    Method & Precision & Recall \\
    \midrule
    Baseline & 0.82 & 0.78 \\
    Proposed & 0.90 & 0.86 \\
    \bottomrule
  \end{tabular}
  \label{tab:results}
\end{table}

\section{Discussion}

Interpret the results, their limitations, and their scope.

\section{Conclusion}

Restate the main contribution and open future perspectives.

\begin{thebibliography}{9}

\bibitem{example}
Author, A.
\newblock Reference title.
\newblock \emph{Journal}, 1(1), 1--10, 2026.

\end{thebibliography}

\end{document}

